\documentclass[12pt,a4paper]{article}
\usepackage[utf8]{inputenc}
\usepackage[T2A]{fontenc}
\usepackage[russian]{babel}
\usepackage{amsmath,amssymb,amsthm}
\usepackage{geometry}
\usepackage{booktabs}
\usepackage{enumitem}

\geometry{left=25mm, right=25mm, top=20mm, bottom=20mm}

\title{Безнапорная фильтрация с запертой газовой шапкой}
\author{}
\date{}

\begin{document}
\maketitle

\section{Постановка задачи}

Рассматривается пласт, ограниченный:
\begin{itemize}[nosep]
  \item снизу~--- непроницаемой подошвой (водоупором),
  \item сверху~--- непроницаемой кровлей.
\end{itemize}
Между ними находится жидкость (снизу) и газовая шапка (сверху), разделённые свободной поверхностью $z = h(x, t)$. Газ \emph{заперт}: он не может уйти через кровлю, его масса фиксирована.

В классической безнапорной фильтрации на свободной поверхности принимается условие $p\big|_{z=h} = p_{\mathrm{atm}} = \mathrm{const}$. В рассматриваемой задаче давление газа является неизвестной, связанной с положением свободной поверхности.

\section{Давление в газовой шапке}

Пусть газ идеальный, процесс изотермический (закон Бойля--Мариотта):
\begin{equation}\label{eq:boyle}
  p_g(t)\, V_g(t) = p_g(0)\, V_g(0) = C = \mathrm{const}.
\end{equation}
Объём газовой шапки определяется конфигурацией свободной поверхности:
\begin{equation}\label{eq:Vg}
  V_g(t) = \int_\Omega \bigl[H(x) - h(x,t)\bigr]\, m(x)\, dx,
\end{equation}
где $H(x)$~--- отметка кровли, $m(x)$~--- пористость, $\Omega$~--- область пласта в плане. Отсюда:
\begin{equation}\label{eq:pg}
  \boxed{p_g(t) = \frac{C}{\displaystyle\int_\Omega \bigl[H(x) - h(x,t)\bigr]\, m(x)\, dx}.}
\end{equation}

Это \emph{нелокальное} (интегральное) соотношение: давление газа в каждый момент времени зависит от \emph{всей} конфигурации свободной поверхности $h(x,t)$ по пласту.

\subsection*{Замечание}
Давление газа $p_g(t)$ одинаково во всех точках газовой шапки~--- газ перераспределяет давление практически мгновенно, поскольку вязкость и плотность газа малы по сравнению с жидкостью. Поэтому $\nabla p_g = 0$.

\section{Граничное условие на свободной поверхности}

Вместо $p = p_{\mathrm{atm}}$ на свободной границе $z = h(x,t)$ имеем:
\begin{equation}\label{eq:bc_free}
  p\big|_{z=h(x,t)} = p_g(t),
\end{equation}
где $p_g(t)$ определяется из~\eqref{eq:pg}.

\section{Уравнение для жидкости (приближение Дюпюи)}

В гидравлическом приближении давление гидростатично по вертикали:
\begin{equation}\label{eq:hydrostatic}
  p(x, z, t) = p_g(t) + \rho g \bigl[h(x,t) - z\bigr].
\end{equation}
Подставляя в закон Дарси и интегрируя по вертикали от подошвы $z=0$ до свободной поверхности $z=h$, получаем:
\begin{equation}\label{eq:bous_full}
  m \frac{\partial h}{\partial t}
  = \nabla \cdot \left[\frac{k h}{\mu}\, \nabla \bigl(p_g(t) + \rho g\, h\bigr)\right] + Q.
\end{equation}
Поскольку $p_g(t)$ не зависит от пространственных координат ($\nabla p_g = 0$), уравнение упрощается:
\begin{equation}\label{eq:bous}
  \boxed{m \frac{\partial h}{\partial t}
  = \nabla \cdot \left[\frac{k \rho g}{\mu}\, h\, \nabla h\right] + Q.}
\end{equation}

Это \emph{то же самое} уравнение Буссинеска, что и в классическом случае. Однако принципиальное отличие состоит в замыкающем соотношении~\eqref{eq:pg} и в модификации граничных условий.

\section{Отличие от классической задачи}

\subsection{Граничные условия с переменным давлением газа}

Если на границе области (скважина, контур питания) задано давление $p_{\mathrm{ext}}$, то напор жидкости у границы:
\begin{equation}\label{eq:h_boundary}
  h\big|_{\partial\Omega} = \frac{p_{\mathrm{ext}} - p_g(t)}{\rho g}.
\end{equation}
Даже при постоянном $p_{\mathrm{ext}}$ эффективный напор $h$ на границе \emph{изменяется во времени} вместе с $p_g(t)$.

\subsection{Газовая пружина}

Запертый газ действует как пружина:
\begin{itemize}[nosep]
  \item подъём $h$ $\to$ уменьшение $V_g$ $\to$ рост $p_g$ $\to$ торможение притока жидкости;
  \item опускание $h$ $\to$ увеличение $V_g$ $\to$ падение $p_g$ $\to$ торможение оттока жидкости.
\end{itemize}
Это создаёт \emph{отрицательную обратную связь}, стабилизирующую систему.

\subsection{Нелокальность}

Откачка жидкости в одной точке пласта изменяет $p_g(t)$, что мгновенно влияет на граничные условия \emph{всюду}. Задача является интегро-дифференциальной.

\subsection{Предельное сжатие}

При $h \to H$ всюду $V_g \to 0$ и $p_g \to \infty$. Физически это означает, что газ не позволяет жидкости полностью заполнить пласт.

\section{Полная постановка задачи}

\begin{equation}\label{eq:full}
  \begin{cases}
    \displaystyle m \frac{\partial h}{\partial t}
    = \nabla \cdot \left[\frac{k \rho g}{\mu}\, h\, \nabla h\right] + Q,
    & x \in \Omega, \quad t > 0, \\[12pt]
    \displaystyle p_g(t) = \frac{C}{\displaystyle\int_\Omega \bigl[H(x) - h(x,t)\bigr]\, m(x)\, dx},
    \\[12pt]
    h\big|_{\partial\Omega} = \dfrac{p_{\mathrm{ext}} - p_g(t)}{\rho g}
    \quad \text{или} \quad
    \left.\dfrac{\partial h}{\partial n}\right|_{\partial\Omega} = \text{задан},
    \\[12pt]
    h(x, 0) = h_0(x).
  \end{cases}
\end{equation}

\section{Обобщения}

\subsection{Адиабатический газ}
Если процесс адиабатический, закон Бойля--Мариотта~\eqref{eq:boyle} заменяется на:
\begin{equation}
  p_g(t)\, V_g^\gamma(t) = p_g(0)\, V_g^\gamma(0),
\end{equation}
где $\gamma$~--- показатель адиабаты. Тогда:
\begin{equation}
  p_g(t) = p_g(0) \left(\frac{V_g(0)}{V_g(t)}\right)^\gamma.
\end{equation}
Нелинейность задачи усиливается ($\gamma > 1$).

\subsection{Учёт капиллярного давления}
Если на границе жидкость--газ существенно капиллярное давление $p_c$, условие~\eqref{eq:bc_free} модифицируется:
\begin{equation}
  p\big|_{z=h} = p_g(t) - p_c,
\end{equation}
и гидростатическое распределение~\eqref{eq:hydrostatic} сдвигается на $p_c$.

\section{Замечания для численной реализации}

\begin{enumerate}[nosep]
  \item На каждом временном шаге необходимо пересчитывать $p_g$ по текущему $h(x,t)$~--- это интеграл по всей области, вычислительно дёшево.
  \item Нелинейность усилена: и уравнение Буссинеска нелинейно, и $p_g(h)$ нелинейна. Может потребоваться итерационная схема (Пикара или Ньютона) на каждом временном шаге.
  \item Связь $p_g$ с $h$ делает задачу \emph{жёсткой} при малых $V_g$ (газ сильно сжат), что может потребовать адаптации шага по времени.
\end{enumerate}

\end{document}
